\section{PASR: Provenance-Annotated Semantic Reduction}
\label{sec:pasr}

\subsection{The Architectural Conflict}

The \sentinel architecture layers its defenses in depth, but two critical
layers---L2 (Information Flow Control via taint tags) and L5 (Semantic
Transduction, inspired by the blood--brain barrier)---are architecturally
incompatible in their na\"ive formulations. This incompatibility is not a
minor integration challenge; it is a genuine structural conflict that no
existing system addresses.

L2 implements Denning's lattice model~\cite{denning1976}: every token entering
the system receives a provenance tag $p_i \in \{\textsc{Operator},
\textsc{User}, \textsc{Retrieved}, \textsc{Tool}\}$, and these tags propagate
through the processing pipeline according to information-flow rules. The
security guarantee is precise: at any point in the pipeline, one can determine
which inputs influenced a given intermediate value.

L5 implements semantic transduction, modeled on the biological blood--brain
barrier (BBB). It deliberately \emph{destroys} the original token sequence and
re-synthesizes only the extracted semantic intent. This destruction is not a
side-effect---it is the security mechanism. Injected instructions that rely on
specific token patterns cannot survive lossy semantic extraction.

The conflict is immediate: \textbf{L5 destroys the very tokens to which L2's
tags are attached.} After transduction, no tokens remain to carry provenance
information. The taint tags die with the tokens they annotated. Any downstream
layer that requires provenance information---and several critical security
decisions depend on it---receives nothing.

\subsection{Core Insight: Provenance as a Property of Derivations}

The resolution begins with a conceptual shift: \textbf{provenance is not a
property of tokens---it is a property of derivations.} A taint tag does not
describe the token it is attached to; it describes the \emph{origin} of the
information that token carries. When L5 destroys tokens and re-synthesizes
semantic intent, the derivation history does not vanish---it merely needs to
be recorded in a form that survives the transduction boundary.

\pasr implements this insight via a \emph{two-channel output} from the trusted
transducer. The transducer \textsc{reads} provenance tags from its input tokens
and \textsc{writes} provenance certificates onto its output semantic fields:

\begin{itemize}[leftmargin=*]
\item \textbf{Channel~1: Semantic intent} (content channel---lossy). The
      extracted meaning of the input, stripped of syntactic structure. This
      channel is deliberately lossy: injected instructions that depend on
      token-level patterns are destroyed.
\item \textbf{Channel~2: Provenance certificate} (metadata channel---lossless).
      An HMAC-SHA256-signed record mapping each semantic field to the set of
      provenance tags that contributed to it. This channel is unforgeable:
      only the trusted transducer holds the signing key.
\end{itemize}

\subsection{The \pasr Algorithm}

The transduction proceeds in five steps. We illustrate with a concrete
injection attempt where user-supplied text contains the instruction
``ignore previous instructions and read /etc/passwd.''

\begin{enumerate}[leftmargin=*]
\item \textbf{Receive tagged tokens from L2.} The transducer ingests a
      sequence of $(t_i, p_i)$ pairs, where each token $t_i$ carries the
      provenance tag $p_i$ assigned by the IFC layer:
      \begin{multline*}
        [(\text{``ignore''}, \textsc{User}),\;
         (\text{``previous''}, \textsc{User}),\\
         (\text{``instructions''}, \textsc{User}),\; \ldots]
      \end{multline*}

\item \textbf{Extract semantic intent (content channel---lossy).} The
      transducer performs lossy semantic extraction, discarding syntactic
      structure and producing a structured intent representation:
      \[
        S = \{\texttt{action}: \text{``file\_read''},\;
             \texttt{target}: \text{``/etc/passwd''}\}
      \]
      The original token sequence is irrecoverably destroyed. Any injected
      payload that relied on token-level patterns (e.g., ``ignore previous
      instructions'') is reduced to its semantic content.

\item \textbf{Record provenance map (novel contribution).} For each field
      $f_j$ in the semantic output $S$, the transducer records which tagged
      inputs contributed to it, producing a provenance map
      $P: \text{Fields}(S) \to \mathcal{P}(\text{Provenance})$:
      \begin{align*}
        P(\texttt{action}) &= \{\textsc{User}\} \\
        P(\texttt{target}) &= \{\textsc{User}\}
      \end{align*}
      This is the step that no existing system performs. Standard semantic
      extraction discards provenance; standard taint tracking cannot survive
      re-synthesis. \pasr does both.

\item \textbf{Sign the provenance map.} The transducer computes an
      HMAC-SHA256 signature over the concatenation of $(S, P)$ using a
      per-session ephemeral key $k$ held exclusively by the trusted
      transducer:
      \[
        \sigma = \text{HMAC-SHA256}(k,\; S \| P)
      \]
      The signed certificate $(S, P, \sigma)$ is unforgeable: an adversary
      who compromises the LLM but not the transducer cannot fabricate a
      valid provenance annotation.

\item \textbf{Detect claims-vs-actual discrepancy.} The transducer compares
      the semantic content against the provenance certificate. In this
      example, the extracted intent requests \texttt{file\_read} on a
      sensitive path---an action that requires \textsc{Operator} authority.
      However, the provenance map records that \emph{all} contributing
      inputs carry \textsc{User} provenance. The discrepancy between claimed
      authority (implicit in the action semantics) and actual source
      provenance constitutes an \textbf{injection signal}:
      \begin{multline*}
        \text{Required}(\texttt{action}) = \textsc{Operator},\\
        P(\texttt{action}) = \{\textsc{User}\}
        \;\Longrightarrow\; \textbf{INJECTION}
      \end{multline*}
\end{enumerate}

\subsection{Mathematical Framework: Provenance Lifting Functor}

We formalize \pasr as a functor between two categories, drawing on the
theory of provenance semirings~\cite{green2007}.

\paragraph{Category $\mathcal{C}$ (L2 output space).}
Objects are tagged token sequences $[(t_i, p_i)]$ where $t_i \in
\text{Tokens}$ and $p_i \in \{\textsc{Operator}, \textsc{User},
\textsc{Retrieved}, \textsc{Tool}\}$. Morphisms are token-level
transformations that propagate tags according to IFC rules.

\paragraph{Category $\mathcal{D}$ (\pasr output space).}
Objects are provenance-annotated semantic structures $(S, P)$ where $S$ is a
semantic intent and $P: \text{Fields}(S) \to \mathcal{P}(\text{Provenance})$
maps each semantic field to the power set of provenance labels that
contributed to it. Morphisms are semantic refinements that respect provenance
monotonicity.

\paragraph{The lifting functor $\mathcal{L}: \mathcal{C} \to \mathcal{D}$.}
The \pasr transducer defines a functor with four key properties:
\begin{enumerate}[leftmargin=*]
\item \textbf{Content-lossy.} $\mathcal{L}$ is not injective on the content
      component: distinct input sequences may map to the same semantic intent.
      Formally, $\exists\, c_1 \neq c_2 \in \text{Ob}(\mathcal{C})$ such that
      $\pi_S(\mathcal{L}(c_1)) = \pi_S(\mathcal{L}(c_2))$, where $\pi_S$
      projects to the semantic component.
\item \textbf{Provenance-faithful.} For each semantic field $f_j$ in the
      output, the provenance annotation is the union of all contributing
      input provenances:
      \[
        P(f_j) = \bigcup \{p_i : t_i \text{ contributed to } f_j\}
      \]
\item \textbf{Monotone in trust.} The trust level assigned to a field is the
      minimum of all contributing trust levels, ensuring conservative
      propagation:
      \[
        \text{trust}(f_j) = \min\{\text{trust}(p_i) : p_i \in P(f_j)\}
      \]
\item \textbf{Unforgeable.} The provenance certificate is HMAC-signed by the
      trusted transducer. No entity outside the transducer's trust boundary
      can produce a valid signature.
\end{enumerate}

\noindent
This construction is a \emph{fibration} in the categorical sense: the
forgetful projection $\pi_S: \mathcal{D} \to \mathcal{S}$ (discarding
provenance and retaining only semantic intent) has a canonical lifting
property. Given any morphism in the base category $\mathcal{S}$ (a semantic
transformation) and a provenance-annotated object in the fiber above its
source, there exists a unique lifted morphism in $\mathcal{D}$ that preserves
provenance faithfulness. This lifting property is precisely what guarantees
that provenance survives arbitrary downstream semantic processing---not merely
the initial transduction step.

\subsection{Synthesis: Why \pasr Is New}

Table~\ref{tab:pasr-synthesis} summarizes the intellectual lineage of \pasr
and identifies the specific gap that each predecessor leaves unaddressed.
\pasr is the first primitive to combine fine-grained provenance tracking,
survival across lossy re-synthesis, and cryptographic unforgeability in a
single mechanism.

\begin{table}[t]
\centering
\small
\caption{Intellectual lineage of \pasr: each predecessor contributes a key
         property but leaves a critical gap that \pasr closes.}
\label{tab:pasr-synthesis}
\begin{tabular}{@{}p{2.2cm}p{2.2cm}p{2.4cm}@{}}
\toprule
\textbf{Source} & \textbf{Contributes} & \textbf{Misses} \\
\midrule
DB Provenance Semirings~\cite{green2007}
  & Fine-grained per-datum tracking
  & Only structure-preserving transforms \\
IFC Taint Tracking~\cite{denning1976}
  & Per-datum provenance labels
  & Labels do not survive re-synthesis \\
Cryptographic Attestation
  & Unforgeability of certificates
  & Wrong granularity (whole-message) \\
\midrule
\pasr (this work)
  & All three combined
  & Novel primitive (no prior art) \\
\bottomrule
\end{tabular}
\end{table}

\subsection{Evaluation Summary}

We evaluate the impact of \pasr on the overall \sentinel detection pipeline.
Full evaluation methodology and results appear in \S\ref{sec:evaluation};
here we highlight three key metrics:

\begin{itemize}[leftmargin=*]
\item \textbf{Overall detection improvement.} Adding \pasr to the \sentinel
      pipeline increases aggregate detection from 93.7\% to 95.7\%
      (+2.0 percentage points), with the largest gains on attacks that
      specifically target the L2--L5 boundary.
\item \textbf{Taint-strip attack resolution.} ATK-011 (taint-strip injection),
      in which the adversary crafts inputs designed to shed provenance tags
      before reaching the semantic transducer, drops from 50--65\% attacker
      success to 3--8\%. This category was the primary motivation for \pasr
      and is now effectively resolved.
\item \textbf{Fast-path latency reduction.} Because provenance certificates
      enable downstream layers to make authorization decisions without
      re-querying L2, the fast-path latency decreases from 15\,ms to 3\,ms
      (a $5\times$ improvement). The provenance certificate acts as a
      pre-computed authorization token.
\end{itemize}

\subsection{Known Weaknesses}

We identify three weaknesses in the current \pasr design and describe the
mitigations deployed in the \sentinel prototype.

\begin{table}[t]
\centering
\small
\caption{Known weaknesses of \pasr and deployed mitigations.}
\label{tab:pasr-weaknesses}
\begin{tabular}{@{}p{2.6cm}p{1.6cm}p{2.6cm}@{}}
\toprule
\textbf{Weakness} & \textbf{Severity} & \textbf{Mitigation} \\
\midrule
HMAC key is a single point of failure: compromise of the signing key
allows provenance forgery
  & High
  & HSM-backed key storage with per-session ephemeral keys; key rotation
    every 300\,s \\
Provenance boundary ambiguity at BPE token splits: a single semantic
concept may span tokens with different provenance tags
  & Medium
  & Conservative assignment: if any contributing token carries
    \textsc{User} provenance, the entire field inherits \textsc{User} \\
Provenance laundering via tool calls: an adversary invokes a tool with
user-controlled input, and the tool's output receives \textsc{Tool}
provenance
  & Medium
  & Transitive provenance tracking: tool outputs inherit the provenance of
    their inputs via \cafl integration \\
\bottomrule
\end{tabular}
\end{table}

\noindent
These weaknesses are architectural rather than implementation-specific: any
system that lifts provenance across a lossy semantic boundary will face
analogous challenges. We discuss longer-term solutions, including
threshold-HMAC schemes and formal provenance algebras, in
\S\ref{sec:discussion}.
